\textbf{First paragraph:}
\begin{itemize}
    \item Context: Diffusion MRI provides quantitative microstructural biomarkers
    \item Increasing use of ML for prediction (age, diagnosis, cognition, etc.)
    \item Lack of standardized benchmark
\end{itemize}

\noindent \textbf{Second paragraph: literature review:}
\begin{itemize}
    \item  Literature review of techniques: Variability of the proposed approach (ML-based, deep-learning based etc)
    \item Problems: High variability in reported performance across studies + Limited reproducibility 
\end{itemize}

% \noindent \textbf{Third paragraph: reasons for these problems}
% \begin{itemize}
%     \item Reason 1: Heterogeneous preprocessing pipelines
%     \item Reason 2: Inconsistent feature representations (voxel-wise vs ROI-based), not analysis cross microstructure features
%     \item Reason 3: Evaluation protocols not directly comparable + results on different datasets
%     \item Reason 4: Missing code
%     \item Our solution: "To address these challenges, we propose: 1) an open-source and standardized benchmark; 2) of various machine learning techniques in dMRI; 3) on several tasks (classification, regression) 4) encompassing N datasets 
% \end{itemize}

These inconsistencies largely stem from four recurrent factors. First, studies often rely on 
heterogeneous preprocessing pipelines (and sometimes different quality-control criteria), 
making results difficult to compare across cohorts. Second, diffusion information is represented in 
markedly different ways (e.g., voxel/vertex-wise maps, tract/ROI summaries, or learned embeddings), 
and evaluations rarely disentangle how the choice of microstructural representation affects downstream 
performance. Third, evaluation protocols vary substantially (e.g., different split strategies, 
cross-validation schemes, and hyperparameter tuning procedures), and models are frequently assessed on 
different datasets, further confounding comparisons. Finally, incomplete release of code, 
configurations, and split definitions limits reproducibility and slows iterative progress. To address 
these challenges, we propose an open-source and standardized benchmark that unifies preprocessing and 
feature extraction, evaluates a diverse set of machine learning approaches on both regression and 
classification tasks, and spans multiple publicly available datasets under a controlled and leakage-free 
evaluation protocol.


% \noindent \textbf{Fourth paragraph: your contributions}
% \begin{itemize}
%     \item Standardized preprocessing pipeline for diffusion MRI
%     \item Implemented various representation of the microstructural input data
%     \item Unified feature extraction strategy
%     \item Controlled evaluation protocol (train/test splits)
%     \item Comparison of multiple feature representations
%     \item Integration of both classical and DL predictors
%     \item Public release of code and benchmark configuration
%     \item End-to-end analysis
% \end{itemize}
In this work, we introduce an open-source, end-to-end benchmark for predictive modeling with 
diffusion MRI that targets reproducibility and fair comparison across studies. Our contributions 
are: (i) a standardized preprocessing and data preparation pipeline; (ii) a unified feature-extraction 
framework supporting multiple microstructural representations for gray and white matter; (iii) a 
controlled evaluation protocol with consistent, leakage-free train/validation/test splits and 
cross-validation; (iv) a systematic comparison of feature representations and predictors, spanning 
classical machine learning baselines and deep-learning models; and (v) the public release of code, 
configurations, and scripts to facilitate transparent reuse and extension.
